\section{The plan on one page}
\label{sec:summary}

\textbf{Goal.} A full-marks attempt at all 40 points: 12 for the report, 20 for the code and zip
file, and 8 HD points. Every task in this plan is tied to one of those points. If a task earns no
point and nothing else depends on it, it was left out.

\textbf{The marks, checked against the brief.} The breakdown used to start this plan missed one
report line. The report is worth 12 points over \emph{six} criteria, not five:

\begin{center}
\small
\begin{tabularx}{\textwidth}{@{}Y C{1.1cm} Y C{1.1cm}@{}}
\toprule
\textbf{Report} (PDF, up to 3,000 words) & \textbf{Pts} & \textbf{Code and zip file} & \textbf{Pts} \\
\midrule
1. Data collection and processing & 5 & 1. Core functionality & 10 \\
2. Data analysis & 2 & 2. Datasets & 4 \\
3. Justification of model selection & 2 & 3. README (conda set-up, training, prediction) & 2 \\
4. Writing quality and logical structure & 1 & 4. Comments & 2 \\
5. Harvard referencing & 1 & 5. Code structure & 2 \\
\textbf{6. Completeness (every required section)} & \textbf{1} & & \\
\midrule
\textit{Report total} & \textit{12} & \textit{Code total} & \textit{20} \\
\midrule
\multicolumn{2}{@{}l}{\textbf{HD}: additional datasets} & 4 & \\
\multicolumn{2}{@{}l}{\textbf{HD}: additional ML implementations} & 4 & \\
\bottomrule
\end{tabularx}
\end{center}

\textbf{What will be built.}
\begin{enumerate}
  \item \textbf{Our own Bitcoin data.} A frozen snapshot of 25 addresses plus 3 short-history ones
    from the Esplora API, as Assignment 1 promised (DS1). A second, much larger crawl of about
    5,000 address histories (DS2) feeds the transaction language model.
  \item \textbf{Two public labelled datasets.} Elliptic (DS3) and BitcoinHeist (DS4) carry labels
    from their authors, so the brief's accuracy, precision, recall and F1 can be reported honestly.
  \item \textbf{All three model types the brief names.} Clustering (K-Means, DBSCAN) and regression
    (Ridge, random forest, gradient boosting) run on our own data. Classification (logistic
    regression, KNN, random forest, gradient boosting) runs on the labelled datasets. The brief
    needs two types; doing all three removes any doubt.
  \item \textbf{The Assignment 1 methods, kept as extra experiments.} The robust-z baseline,
    Isolation Forest and Local Outlier Factor are still built and evaluated with the label-free
    plan. They are also tested against the Elliptic labels, which shows how often ``unusual''
    and ``labelled illicit'' coincide.
  \item \textbf{Two language models trained from scratch on the RTX 3060.} The first, TxGPT, is a
    small GPT that reads an address's history as a sequence of transaction tokens. It scores how
    surprising each new transaction is, and it is the main HD model. The second is a 152M-parameter
    text model, pretrained on 2B tokens of web text and then taught to chat and to explain flags.
    It becomes the Explorer Assistant, which Assignment 3 can use.
  \item \textbf{A prediction interface.} \code{predict.py} plus a command-line tool, so every model
    can score new input. The README shows how, and the Assignment 3 FastAPI back end can call
    the same functions.
\end{enumerate}

\textbf{How this keeps the Assignment 1 promise.} Assignment 1 ruled out claiming fraud accuracy on
data with no labels (risk 13). That rule still holds. Precision and recall are reported only in two
places. One is datasets whose authors supplied labels. The other is transactions we planted
ourselves, where the true answer is known by construction. The Explorer's wording never changes:
\emph{``Unusual activity requires investigation; it does not indicate wrongdoing.''}

\textbf{The dates that matter.}
\begin{center}
\small
\begin{tabularx}{\textwidth}{@{}L{3.1cm}Y@{}}
\toprule
\textbf{Date} & \textbf{Milestone} \\
\midrule
Wed 21 Oct & Change request CR-06 sent to the team and the tutor. Snapshot DS1 fetched. Text pretraining started on the GPU. \\
Sun 25 Oct & M1: baseline, Isolation Forest and LOF score every snapshot transaction. This is Assignment 1 Objective 2, due 26 Oct. \\
Thu 29 Oct & M2: every model trained and evaluated; results frozen (version 1). \\
Sat 31 Oct & M3: code feature-complete; notebooks re-run from clean. \\
Wed 4 Nov & M4: a teammate rebuilds everything from the README on a clean machine; code frozen. \\
Sat 7 Nov & M5: everything submitted on Canvas. Sunday 8 Nov is spare. \\
\bottomrule
\end{tabularx}
\end{center}

\textbf{The critical path.} DS1 fetch → features → time split and baseline → models → evaluation
→ method selection → report evaluation section → submission. Both language models run on a
parallel GPU lane with about three days of slack. If the text model's pretraining runs long, a
decision gate cuts its token budget (Section~\ref{sec:gate}). It never delays the critical path.
